% Proposal — submitted Nov 10, 2 pages, 11 pt.
% Rewritten to use the shared deqreport.sty so all four documents match.
% Build: make proposal
\documentclass[11pt]{article}
\usepackage{deqreport}
\reporttitle{Final Project Proposal}

\begin{document}
\makeheader

\section{Topic}

This project follows the Efficient Implementation track, Direction~2: I will
implement a deep equilibrium model (DEQ) myself, including the fixed-point
forward solver and the implicit-differentiation backward pass. Automatic
differentiation will not be run through the solver iterations.

The deliverable has two parts. The first is a small fully connected DEQ in
NumPy with every gradient derived and coded by hand, verified by
finite-difference gradient checks. The second is a convolutional DEQ in
PyTorch whose backward pass is a custom function I write, trained on MNIST and
CIFAR-10. I will use it to compare three ways of computing the gradient (exact
implicit, truncated Neumann, and Jacobian-free) against an unrolled weight-tied
baseline on accuracy, memory, and run time.

\section{Motivation}

A DEQ replaces a stack of $L$ layers with one layer applied until its output
stops changing, so the network is defined by an equation instead of a fixed
sequence of operations. For input $x$, the hidden state is the fixed point
\begin{equation}
  \zstar = \fth(\zstar, x).
\end{equation}
Because the gradient depends only on the equilibrium and not on the path the
solver took, training needs constant memory in the effective depth:
\begin{equation}
  \dd{\ell}{\theta}
  = \dd{\ell}{\zstar}\left(I - \dd{\fth}{\zstar}\right)^{-1}\dd{\fth}{\theta}.
  \label{eq:implicit}
\end{equation}
This is attractive, but it moves the difficulty elsewhere. Standard
backpropagation, the rule taught in this course, no longer applies directly:
the backward pass becomes a linear solve, and the forward pass may converge
slowly or not at all as the weights grow during training. Frameworks hide these
issues behind library calls. Implementing the solver and the gradient by hand
is the most direct way to understand them, and it allows a controlled answer to
a practical question: how much accuracy is lost when the exact inverse in
Eq.~\eqref{eq:implicit} is replaced by a cheap approximation, and how much time
is saved?

\section{Related Work}

\begin{itemize}
  \item \textbf{Implicit layers.} OptNet~\cite{amos2017optnet} differentiates
    through a quadratic program, Neural ODEs~\cite{chen2018neural} through an
    ODE solver using the adjoint method, and El~Ghaoui et
    al.~\cite{elghaoui2021implicit} give a general formulation and
    well-posedness conditions for implicit models.
  \item \textbf{Deep equilibrium models.} Bai et al.~\cite{bai2019deq}
    introduce the DEQ for sequence modeling, using Broyden's
    method~\cite{broyden1965class} for both passes. Multiscale
    DEQ~\cite{bai2020mdeq} extends it to image classification and segmentation.
  \item \textbf{Convergence guarantees.} Monotone operator equilibrium
    networks~\cite{winston2020monotone} parameterize the layer so that a unique
    fixed point exists and operator-splitting solvers are guaranteed to find it.
  \item \textbf{Cheaper and more stable training.} Jacobian
    regularization~\cite{bai2021stabilizing} penalizes the Jacobian norm at
    equilibrium to keep solver iterations from growing. Phantom
    gradients~\cite{geng2021training} and Jacobian-free
    backpropagation~\cite{fung2022jfb} replace the exact inverse with a
    truncated or zeroth-order approximation.
\end{itemize}

The gap this project addresses is a small, fully hand-derived implementation in
which the three backward methods are compared under one architecture, one
solver, and one compute budget.

\section{Potential Solutions}

\begin{enumerate}
  \item \textbf{Layer.} The fully connected version uses
    $f(z, x) = \ReLU(Wz + Ux + b)$. The convolutional version uses two
    $3\times 3$ convolutions with group normalization and input injection. $W$
    is rescaled by its spectral norm so the layer starts as a contraction.
  \item \textbf{Forward solver.} Plain fixed-point iteration is the baseline.
    Anderson acceleration~\cite{anderson1965iterative} with a memory of 5 is
    the main solver, stopping at a relative residual of $10^{-3}$ or 50
    iterations.
  \item \textbf{Backward pass.} Three variants share the same forward pass:
    \begin{itemize}
      \item Exact implicit gradient: solve the linear system
        $u = uJ + \dd{\ell}{\zstar}$ with the same Anderson solver.
      \item Truncated Neumann series: approximate the inverse by its first $k$
        terms, for $k$ from 1 to 5.
      \item Jacobian-free: replace the inverse with the identity ($k = 0$).
    \end{itemize}
  \item \textbf{Baseline.} A weight-tied network that applies the same layer
    $K$ times and trains with ordinary backpropagation through all $K$ steps.
\end{enumerate}

\textbf{Evaluation.} I will report test accuracy, peak memory, time per epoch,
forward and backward iteration counts over training, and the cosine similarity
between each approximate gradient and the exact one.

\textbf{Risks.} If the forward solver stops converging late in training, I will
add Jacobian regularization~\cite{bai2021stabilizing}. If CIFAR-10 exceeds a
single-GPU budget, the full comparison will run on MNIST and CIFAR-10 will use
a reduced model.

\section{Timeline}

\textbf{Nov 17 (progress report 1):} NumPy DEQ on MNIST with fixed-point and
Anderson solvers, exact implicit backward pass, gradient check passing.
\textbf{Dec 1 (progress report 2):} convolutional DEQ on CIFAR-10, Neumann and
Jacobian-free backward passes, unrolled baseline, first comparison table.
\textbf{Dec 8 (final report):} ablations on solver tolerance and Anderson
memory, final results, write-up.

\printrefs
\end{document}
